%TCIDATA{Version=5.50.0.2960}
%TCIDATA{LaTeXparent=0,0,main.tex}
                      

\pretolerance=2000 \tolerance=3000

\nonumchapter{Lista de s�mbolos}

\begin{tabular}{ll}
$\mathbf{R}$ & Conjunto de todos los n\'{u}meros reales \\ 
$\mathbf{R}^{n}$ & Espacio $n-$dimensional, $\mathbf{R}^{n}=\mathbf{R}\times
\cdots \times \mathbf{R}$ $\left( n\text{ veces}\right) $ \\ 
$\mathbf{R}_{+,0}^{n}$ & Ortante no negativo, subconjunto de n\'{u}meros no
negativos de $\mathbf{R}^{n}$ \\ 
$\mathbf{R}^{n\times n}$ & Indica una matriz cuadrada de $n\times n$
elementos \\ 
$R_{0}$ & N\'{u}mero reproductivo b\'{a}sico \\ 
$R^{2}$ & Coeficiente de determinaci\'{o}n \\ 
$P$ & Indica que la din\'{a}mica de un sistema se localiza en $\mathbf{R}%
_{+,0}^{n}$ \\ 
$\overset{.}{x}$ & La primera derivada con respecto al tiempo de $x$, notaci%
\'{o}n de Newton \\ 
$dx/dt$ & La primera derivada con respecto al tiempo de $x$, notaci\'{o}n de
Leibniz \\ 
$f\left( x\right) $, $f\left( t,x\right) $ & Vector de estados que describe
el modelo matem\'{a}tico de un sistema din\'{a}mico \\ 
$\overset{.}{x}=f\left( t,x\right) $ & Sistema din\'{a}mico no lineal,
continuo y variante en el tiempo \\ 
$\overset{.}{x}=f\left( x\right) $ & Sistema din\'{a}mico no lineal,
continuo e invariante en el tiempo \\ 
$\overset{.}{x}=Ax+b$ & Sistema din\'{a}mico lineal, continuo e invariante
en el tiempo \\ 
$\phi \left( t,x\right) $ & Indica la soluci\'{o}n en el tiempo de $\overset{%
.}{x}$ \\ 
$x\left( t\right) $ & Indica una soluci\'{o}n en series de tiempo o la soluci%
\'{o}n an\'{a}litica de $\dot{x}$ \\ 
$x\left( 0\right) =x_{0}$ & Condici\'{o}n inicial \\ 
$\phi \left( t,x_{0}\right) $ & Soluci\'{o}n de $\overset{.}{x}$ con la
condici\'{o}n inicial $x_{0}$ \\ 
$\phi \left( \cdot ,x\right) $ & Trayectoria en el espacio de fase
correspondiente a las soluciones de $\dot{x}$ \\ 
$\left\{ {}\right\} $ & Se escribe para indicar los elementos de un
conjunto, $U=\left\{ x\text{ }|\text{ }x\in \mathbf{R}^{n}\text{ }\right\} $
\\ 
$\emptyset $ & Conjunto vac\'{\i}o \\ 
$\partial U$ & Frontera del conjunto $U$ \\ 
$\left\{ x_{i}\right\} $ & Secuencia, el orden de los elementos $x_{i}$ del
conjunto es relevante \\ 
$\subset $ & Subconjunto \\ 
$\cup $ & Uni\'{o}n \\ 
$\cap $ & Intersecci\'{o}n \\ 
$\in $ & Pertenece a \\ 
$\forall $ & Para todo \\ 
$|$, $:$ & Tal que \\ 
$\rightarrow $ & Tiende a \\ 
$\Longrightarrow $ & Implica que \\ 
$\Longleftrightarrow $ & Equivalente a, si y solo si \\ 
$f:U\rightarrow R$ & La funci\'{o}n $f$ mapea el conjunto $U$ al conjunto $R$%
\end{tabular}

\begin{tabular}{ll}
$\sup $ & Valor supremo, la menor de las cotas superiores \\ 
$\inf $ & Valor \'{\i}nfimo, la mayor de las cotas inferiores \\ 
$\Gamma $ & Indica una trayectoria \\ 
$\Gamma _{x}^{+}$ & Media trayectoria (semi trayectoria) positiva, para $%
t\geq 0$ \\ 
$\Gamma _{x}^{-}$ & Media trayectoria (semi trayectoria) negativa, para $%
t\leq 0$ \\ 
$\omega \left( \Gamma \right) $ & Conjunto $\omega -$l\'{\i}mite \\ 
$\alpha \left( \Gamma \right) $ & Conjunto $\alpha -$l\'{\i}mite \\ 
$N(x,\varepsilon )$ & Vecindad \\ 
$\lim $ & L\'{\i}mite \\ 
$\det $ & Determinante \\ 
$:=$ & Igual por definici\'{o}n \\ 
$\equiv $ & Equivalente a \\ 
$\exists $ & Existe \\ 
$T$ & Transpuesta de un vector o periodo de una \'{o}rbita peri\'{o}dica \\ 
$\Delta _{x}$ & $\Delta _{x}=x_{i+1}-x_{i}$ \\ 
$\square $ & Indica el final de una prueba \\ 
$\left\Vert x\right\Vert $ & Norma de $x$ \\ 
$d\left( x,y\right) =\left\vert \left\vert x-y\right\vert \right\vert $ & 
Distancia entre los vectores $x$ y $y$ \\ 
$\dfrac{\partial f\left( x\right) }{\partial x}$ & Matriz Jacobiana \\ 
$\lambda _{i}$ & Valores propios de la matriz Jacobiana o exponentes de
Lyapunov \\ 
$\func{Re}$ & Parte real de un n\'{u}mero complejo \\ 
$\func{Re}\lambda _{i}$ & Parte real del valor propio \\ 
$C^{1}$ & Indica que una funci\'{o}n es de clase continuamente diferenciable
\\ 
$C^{\infty }$ & Indica que una funci\'{o}n es de clase infinitamente
diferenciable \\ 
$\Psi $ & Conjunto compacto invariante \\ 
$h\left( x\right) $ & Notaci\'{o}n empleada para escribir una funci\'{o}n
localizadora \\ 
$L_{f}h\left( x\right) $ & Derivada de Lie de la funci\'{o}n $h\left(
x\right) $ \\ 
$S(h)$ & Notaci\'{o}n empleada para escribir el conjunto $\{x\in \mathbf{R}%
^{n}\mid L_{f}h(x)=0\}$ \\ 
$x^{\ast }$ & Punto de equilibrio \\ 
$V\left( x\right) $ & Funci\'{o}n candidata de Lyapunov \\ 
$\dot{V}\left( x\right) $ & Derivada de la funci\'{o}n candidata de Lyapunov
\\ 
$\alpha $ & Nivel de significaci\'{o}n \\ 
$dof$ & Grados de libertad \\ 
$\Sigma $ & Matriz de covarianzas \\ 
$SE$ & Error estand\'{a}r \\ 
$CI$ & Intervalos de confianza%
\end{tabular}
